% Elemento opcional, redigido e aprovado pelo autor.
% Exemplo do modelo, substituído pelo autor: Agradeço primeiramente...
\begin{agradecimentos}

A conclusão deste trabalho representa não apenas o encerramento de uma etapa acadêmica, mas também o resultado de uma trajetória construída entre estudos, responsabilidades profissionais, desafios pessoais e o desejo contínuo de aprender.

Agradeço ao Prof. Dr. Rafael Zanin, pela orientação, pela leitura criteriosa e pelas contribuições fundamentais para o desenvolvimento deste trabalho. Ao Prof. Dr. Manassés Ribeiro, pela coorientação, pela disponibilidade e pelas reflexões que contribuíram para o amadurecimento da pesquisa.

Aos professores do curso de Bacharelado em Ciência da Computação do Instituto Federal Catarinense -- Campus Videira, pelos conhecimentos compartilhados e pela contribuição à minha formação acadêmica e profissional. Aos membros da banca, Prof. Me. Fábio José Rodrigues Pinheiro e Prof. Dr. Dani Prestini, pelo tempo dedicado à avaliação deste trabalho.

À minha mãe, que foi e sempre será minha base, por acreditar em mim e incentivar meus estudos ao longo da vida. Seu apoio constante e sua confiança na importância da educação fazem parte desta conquista.

À minha família, pelo apoio, pela compreensão e por fazer parte de uma trajetória em que nem sempre foi simples conciliar as responsabilidades da vida com as exigências da formação acadêmica. A todos que estiveram ao meu lado, oferecendo incentivo, companhia e compreensão ao longo desse percurso, minha gratidão.

Ao meu filho, Murilo, que ocupa um lugar especial na minha vida e representa uma das maiores motivações para continuar aprendendo e construindo novos caminhos. Que esta conquista também lhe mostre que o conhecimento, a curiosidade e a perseverança podem nos levar além daquilo que inicialmente imaginamos ser possível.

Por fim, agradeço ao Instituto Federal Catarinense e à educação pública, pelas oportunidades de formação e pelo papel que desempenham na construção de uma sociedade em que o conhecimento possa transformar realidades.

\end{agradecimentos}
